% ch3_c 第3章 §3.3–§3.4 (书页 90–100 / p105–p115)
% 边界备注：起点——书页 90 顶部为完整新段（书页 89 止于 3.3.7 节要点列表），
% 无接缝标记（JOIN）、无上一块收尾（TAIL）可写入；
% 终点——书页 100 止于完整段落（"……描写了超流转变。"），书页 101 顶起新段
% "In the superfluid phase, ..."，本块自然收束，无 TAIL。

在 3.3.2 节中我们提到，超流相（或对称性破缺相）由非零序参量 $\varphi_0$ 表征。在量子理论中，序参量是玻色子场在基态上的期待值 $\varphi_0=\langle\Phi_0|\varphi(\pmb{x})|\Phi_0\rangle=\langle\Phi_0|a(\pmb{x})|\Phi_0\rangle$（记住我们已把玻色子湮灭算符 $a(\pmb{x})$ 改记为 $\varphi(\pmb{x})$）。对于一个有限的量子系统，序参量恒为零，不存在自发对称性破缺。这是因为在量子层次上，$U(1)$ 对称性由算符 $W=\mathrm{e}^{-\mathrm{i}\hat{N}\theta}$ 生成，其中 $\hat{N}$ 是总玻色子数算符。容易验证
\begin{equation*}
Wa(x)W^{\dagger}=\mathrm{e}^{\mathrm{i}\theta}a(x) \tag{3.3.19}
\end{equation*}
由于 $[H,W]=0$，有限系统的基态 $|\Phi_0\rangle$ 总是 $W$ 的本征态。于是 $\langle\Phi_0|a(\pmb{x})|\Phi_0\rangle=\langle\Phi_0|Wa(\pmb{x})W^{\dagger}|\Phi_0\rangle$，因而序参量 $\langle\Phi_0|a(\pmb{x})|\Phi_0\rangle=\langle\Phi_0|\varphi(\pmb{x})|\Phi_0\rangle=0$，无论 $\mu$ 取何值。

与之相对，3.3.2 节中的经典计算确实预言了相变和破缺 $U(1)$ 对称性的简并基态，即使对有限系统也是如此。（事实上，上一节的所有计算都是对体积为 $V$ 的有限系统进行的。）$U(1)$ 对称性之所以能在有限的经典系统中破缺，是因为 $\varphi$ 场的均匀部分 $\varphi_0$ 没有涨落。下面我们将看到，如果把 $\varphi_0$ 当作量子变量处理，它的量子涨落将恢复有限系统的 $U(1)$ 对称性。我们还将看到，当系统体积 $\mathcal{V}\to\infty$ 时，$\varphi_0$ 的量子涨落趋于零。因此，$U(1)$ 对称性破缺可以在量子系统的 $\mathcal{V}\to\infty$ 极限下演生。

我们知道，低能物理由 XY 模型支配。这里我们只关心均匀涨落，于是在式 (3.3.10) 中取 $\theta(\pmb{x},t)=\theta_0(t)$。支配 $\theta_0$ 动力学的有效理论由下式给出
\begin{equation*}
L=\mathcal{V}\big[\frac{1}{2V_0}(\partial_t\theta_0)^2-\rho_0\partial_t\theta_0+\mu\rho_0-\frac{1}{2}V_0\rho_0^2\big] \tag{3.3.20}
\end{equation*}
其中我们放回了一个全时间导数项和一些常数项。式 (3.3.20) 描写一个简单的系统：一个在由 $\theta_0\in[0,2\pi]$ 参数化的圆上运动的粒子。粒子的质量为 $M=\mathcal{V}/V_0$，动量为
\begin{equation*}
p=\partial_{\dot{\theta}_0}L=M\dot{\theta}_0-\frac{\mu\mathcal{V}}{V_0} \tag{3.3.21}
\end{equation*}
其中用到了 $\mu=\rho_0V_0$。哈密顿量为
\begin{equation*}
H=p\dot{\theta}_0-L=\frac{p^2}{2M}+\frac{\mu\mathcal{V}p}{V_0}
\end{equation*}

如果把 $H$ 当作经典系统处理，那么当粒子具有动量 $p=-\mu\mathcal{V}/V_0$（译注：原文为 $-\mu\mathcal{V}p/V_0$，疑多排了一个 $p$），即速度 $\dot{\theta}=0$ 时，它处于基态。粒子的位置可以任意。因此，描写粒子处于不同 $\theta_0$ 的所有态 $|\theta_0\rangle$ 都是简并的基态，$U(1)$ 对称性破缺。

然而，量子理论给出了一幅截然不同的图像。在量子理论中，动量 $p$ 按整数量子化。动量为 $p=-n$ 的态具有能量
\begin{equation*}
E_n=\frac{n^2}{2M}-\frac{\mu\mathcal{V}n}{V_0M}=\frac{V_0n^2}{2\mathcal{V}}-\mu n \tag{3.3.22}
\end{equation*}
因此，$-p=n=\mu\mathcal{V}/V_0$ 的态能量最小，对应于真正的基态
\begin{equation*}
|\Phi_0\rangle=\int\mathrm{d}\theta_0\,\mathrm{e}^{-\mathrm{i}\frac{\mu\mathcal{V}}{V_0}\theta_0}|\theta_0\rangle
\end{equation*}
基态没有简并。现在可以明确地证明序参量为零：
\begin{equation*}
\langle\Phi_0|\varphi_0\mathrm{e}^{\mathrm{i}\theta_0}|\Phi_0\rangle=0
\end{equation*}

理解为什么没有 $U(1)$ 对称性破缺的另一种方式，是注意到总玻色子数为 $N=\mathcal{V}(\rho_0-\frac{\partial_t\theta_0}{V_0})$（见式 (3.3.13)）。由式 (3.3.21) 可见，$-N=p$ 正是 $\theta_0$ 的正则动量（canonical momentum）。因此，在量子世界中二者之间存在不确定关系（uncertainty relation）：
\begin{equation*}
[\hat{N},\theta_0]=\mathrm{i},\qquad \Delta N\,\Delta\theta_0\sim 1/2
\end{equation*}
所以，粒子数固定的有限系统不可能具有确定的相位 $\theta_0$。要得到对称性破缺态，我们必须允许粒子数容易地涨落。\footnote{因此，如果添加或移走一个玻色子需要花费有限的能量，那么该系统就不可能处于 $U(1)$ 对称性破缺相。作为这一理解的一个应用，让我们考虑一个具有强在位排斥（on-site repulsion）的格点玻色子系统 $H=\sum_{\langle ij\rangle}t_{ij}a_i^{\dagger}a_j+\sum_r U(a_r^{\dagger}a_r)^2$，其中 $U\gg t_{ij}$。若玻色子密度恰好是每格点整数个玻色子，则该系统不可能处于超流相，因为在这个密度下添加或移走一个玻色子需要花费有限的能量。}

虽然量子涨落消除了简并、恢复了有限系统的 $U(1)$ 对称性，但对于大系统，低激发态是近乎简并的。能隙只有 $V_0/\mathcal{V}$ 的量级，当 $\mathcal{V}\to\infty$ 时趋于 0。如果该能隙低于所有我们感兴趣的能标（例如低于测量的能量分辨率），那么在一切实际用途中，这些低激发态都可以认为是简并的。这时，有限系统可以认为具有 $U(1)$ 对称性破缺。在问题 3.3.4 中我们将看到，如果我们制备一个相位为 $\theta_0$ 的态，那么当系统很大时，相位要经过很长时间才能变到其他值。因此，在短的时间间隔内，我们可以把相位当作刻画对称性破缺基态的一个常数。

在 3.3.2 节的末尾我们讨论过序参量和长程关联。本节讨论一类特殊的量子涨落——$\varphi$ 的均匀相位涨落。我们看到，对于有限系统，即使 $\mu>0$，这类涨落也会破坏序参量。但同样清楚的是，$\mu>0$ 时存在的长程关联不会被均匀相位涨落破坏。因此，对于大的有限系统，我们可以用长程关联来刻画对称性破缺相变。

在上面的讨论中，我们只考虑了 $\theta$ 的均匀涨落，忽略了依赖空间的涨落。这里的问题是，什么时候这样做才是正确的？我们注意到，依赖空间的涨落对应于具有有限动量的声波。这些声波激发的最小能隙为 $2\pi v/L$，其中 $L$ 是系统的线度。于是，在能量 $2\pi v/L$ 以下，我们可以忽略依赖空间的 $\theta$ 涨落，只考虑均匀涨落。本节的讨论仅在 $2\pi v/L$ 以下有效。

具有固定角度的对称性破缺态 $|\theta_0\rangle$ 由 $|\theta_0\rangle=\sum_n\mathrm{e}^{\mathrm{i}\theta_0 n}|n\rangle$ 给出。要在上述近似之内构造这样一个态，我们需要在能量 $2\pi v/L$ 以下容纳无穷多个均匀涨落。在一维以外，均匀涨落的能隙为 $V_0/\mathcal{V}$，远小于声波的能隙。在 $L\to\infty$ 极限下，$2\pi v/L$ 以下有无穷多条均匀涨落的能级。因此，在一维以外可以构造对称性破缺态。在一维，$2\pi v/L$ 以下的均匀态数目有限，我们难以构造出对称性破缺态。在 3.3.8 节中我们将看到，一维的超流相不可能具有真正的长程序。

\subsection*{问题 3.3.4}

考虑零温下 $1\,\mathrm{cm}^3$ 的超流 $\mathrm{He}_4$。我们制备一个相位 $\langle\theta_0\rangle=0$ 的“基态”。假设相位的弥散为 $\sqrt{\langle\theta_0^2\rangle}=\pi/10$。问经过多长时间相位的弥散会达到 $2\pi$，使得基态的相位不再能定义？（提示：你需要猜测 $\mathrm{He}_4$ 中声波速度的数量级，以便估计 $V_0$——$V_0$ 恰好是压缩率（compressibility）的倒数。）

\subsection{低维中的超流相}

\begin{keypoints}
\item 与 Nambu--Goldstone 模相联系的量子涨落可以破坏低维中的长程序和超流相。
\item 量子涨落对量子临界点（quantum critical point）的影响，以及上临界维数（upper critical dimension）的概念。
\end{keypoints}

在 3.3.2 节的末尾，我们证明了对称性破缺相中长程序的存在。然而，那个计算几乎是在作弊，因为只考虑经典基态，就不允许存在依赖空间的相位涨落，其结果是不同位置处的相位总是关联的。这里我们将研究依赖空间的相位涨落的效应，看看长程关联何时能在依赖空间的相位涨落下存活。

首先考虑零温下的玻色子系统。我们想计算 $\langle\varphi^{\dagger}(t,\pmb{x})\varphi(0,0)\rangle$。由于只关心低能激发，我们将使用有效 XY 模型（式 (3.3.10)），它可以改写为
\begin{equation*}
L=\frac{\chi}{2}\big((\partial_t\theta)^2-v^2(\partial_{\pmb{x}}\theta)^2\big) \tag{3.3.23}
\end{equation*}
其中 $\chi=1/V_0$ 是超流的压缩率。在路径积分方法中，关联函数可以写为
\begin{equation*}
\big\langle\varphi(t,\pmb{x})^{\dagger}\varphi(0,0)\big\rangle=|\varphi_0|^2\big\langle\mathrm{e}^{-\mathrm{i}\theta(t,\pmb{x})}\mathrm{e}^{\mathrm{i}\theta(0,0)}\big\rangle=\frac{\int\mathcal{D}\theta\,\mathrm{e}^{-\mathrm{i}\theta(t,\pmb{x})}\mathrm{e}^{\mathrm{i}\theta(0,0)}\mathrm{e}^{\mathrm{i}S}}{\int\mathcal{D}\theta\,\mathrm{e}^{\mathrm{i}S}}
\end{equation*}
引入
\begin{equation*}
Z[f(t,\pmb{x})]=\int\mathcal{D}\theta\,\mathrm{e}^{\mathrm{i}S}\mathrm{e}^{\mathrm{i}\int\mathrm{d}t\,\mathrm{d}^d\pmb{x}\,f(t,\pmb{x})\theta(t,\pmb{x})}
\end{equation*}
我们看到
\begin{equation*}
\big\langle\mathrm{e}^{-\mathrm{i}\theta(t_0,\pmb{x}_0)}\mathrm{e}^{\mathrm{i}\theta(0,0)}\big\rangle=\frac{Z[-\delta(t-t_0,\pmb{x}-\pmb{x}_0)+\delta(t,\pmb{x})]}{Z[0]}
\end{equation*}
这里 $Z[f(t,\pmb{x})]/Z[0]$ 可以用高斯积分容易地算出：
\begin{equation*}
\frac{Z[f(t,\pmb{x})]}{Z[0]}=\mathrm{e}^{-\mathrm{i}\frac{1}{2}\int\mathrm{d}t_1\mathrm{d}\pmb{x}_1\mathrm{d}t_2\mathrm{d}\pmb{x}_2\,f(t_1,\pmb{x}_1)G_{\theta}(t_1-t_2,\pmb{x}_1-\pmb{x}_2)f(t_2,\pmb{x}_2)}
\end{equation*}
其中 $G_{\theta}$ 是 $\chi(-\partial_t^2+v^2\partial_{\pmb{x}}^2)$ 的逆。它也是 $\theta$ 场的关联函数：
\begin{equation*}
\mathrm{i}G_{\theta}(t,\pmb{x})=\langle T[\theta(t,\pmb{x})\theta(0,0)]\rangle
\end{equation*}
我们发现
\begin{equation*}
\big\langle\mathrm{e}^{-\mathrm{i}\theta(t,\pmb{x})}\mathrm{e}^{\mathrm{i}\theta(0,0)}\big\rangle=\mathrm{e}^{\mathrm{i}G_{\theta}(t,\pmb{x})}\mathrm{e}^{-\mathrm{i}G_{\theta}(0,0)}\propto\mathrm{e}^{\langle(-\mathrm{i}\theta(t,\pmb{x}))(\mathrm{i}\theta(0,0))\rangle}
\end{equation*}

\subsubsection{$1+1$ 维中的超流相}

$k$--$\omega$ 空间中的关联函数由下式给出
\begin{equation*}
G_{\theta}(\omega,\pmb{k})=\frac{\chi^{-1}}{\omega^2-v^2\pmb{k}^2+\mathrm{i}0^+}.
\end{equation*}
我们可以用它按如下方式计算 $1+1$ 维的 $G_{\theta}(t,\pmb{x})$：
\begin{align*}
G_{\theta}(t,x)&=\int\frac{\mathrm{d}k\,\mathrm{d}\omega}{(2\pi)^2}\,G_{\theta}(\omega,k)\,\mathrm{e}^{\mathrm{i}(-\omega t+kx)}\\
&=\chi^{-1}\int\frac{\mathrm{d}k\,\mathrm{d}\omega}{(2\pi)^2}\,\frac{1}{2v|k|}\left(\frac{1}{\omega-v|k|+\mathrm{i}0^+}-\frac{1}{\omega+v|k|-\mathrm{i}0^+}\right)\mathrm{e}^{\mathrm{i}(\omega t-kx)}\\
&=\left\{\begin{array}{ll}
\chi^{-1}\displaystyle\int\frac{\mathrm{d}k}{(2\pi)^2}\,\frac{-\mathrm{i}\pi}{v|k|}\,\mathrm{e}^{\mathrm{i}(-v|k|t+kx)}, & t>0\\[10pt]
\chi^{-1}\displaystyle\int\frac{\mathrm{d}k}{(2\pi)^2}\,\frac{-\mathrm{i}\pi}{v|k|}\,\mathrm{e}^{\mathrm{i}(+v|k|t+kx)}, & t<0
\end{array}\right.\\
&=\chi^{-1}\int\frac{\mathrm{d}k}{(2\pi)^2}\,\frac{-\mathrm{i}\pi}{v|k|}\,\mathrm{e}^{\mathrm{i}(-v|k||t|+kx)}
\end{align*}
对于有限系统，$k$ 是量子化的：$k=\frac{2\pi}{L}\times\text{整数}$。把 $\int\mathrm{d}k/2\pi$ 换成 $L^{-1}\sum_k$ 后，我们得到
\begin{align*}
G_{\theta}(t,x)&=\chi^{-1}L^{-1}\sum_k\frac{-\mathrm{i}}{2v|k|}\,\mathrm{e}^{\mathrm{i}(-v|k||t|+kx)}\\
&=\frac{\mathrm{i}}{4\pi v\chi}\left(\ln\big(1-\mathrm{e}^{2\pi\mathrm{i}\frac{-v|t|+x}{L}-0^+}\big)+\ln\big(1-\mathrm{e}^{2\pi\mathrm{i}\frac{-v|t|-x}{L}-0^+}\big)\right)
\end{align*}
两个对数项分别来自 $k>0$ 与 $k<0$ 的求和。$k=0$ 项是常数，已被略去。如果 $v|t|,\,x\ll L$，则
\begin{equation*}
G_{\theta}(t,x)=\frac{\mathrm{i}}{4\pi v\chi}\ln 4\pi^2\frac{x^2-v^2t^2+\mathrm{i}0^+}{L^2} \tag{3.3.24}
\end{equation*}
我们看到，在 $1+1$ 维中，
\begin{align*}
\big\langle\mathrm{e}^{-\mathrm{i}\theta(t,x)}\mathrm{e}^{\mathrm{i}\theta(0,0)}\big\rangle&=\mathrm{e}^{-\mathrm{i}G_{\theta}(0,0)}\left(\frac{L^2}{4\pi^2(x^2-v^2t^2+\mathrm{i}0^+)}\right)^{1/4\pi v\chi}\\
&=\mathrm{e}^{-\mathrm{i}G_{\theta}(0,0)}\,\mathrm{e}^{-\mathrm{i}\frac{\pi}{4\pi v\chi}\Theta(v^2t^2-x^2)}\left(\frac{L^2}{4\pi^2|x^2-v^2t^2|}\right)^{1/4\pi v\chi}
\end{align*}
由式 (3.3.24)，$G_{\theta}(0,0)=-\mathrm{i}\infty$，这似乎意味着 $\langle\mathrm{e}^{-\mathrm{i}\theta(t,x)}\mathrm{e}^{\mathrm{i}\theta(0,0)}\rangle=0$。我们提醒，XY 模型只是一个低能有效理论，它在短距离处失效。$G_{\theta}(0,0)$ 的发散应当被一个短距离标度 $l$ 截断。于是，我们可以把 $G_{\theta}(0,0)$ 换成 $G_{\theta}(0,l)$，这就得到
\begin{align*}
\big\langle\mathrm{e}^{-\mathrm{i}\theta(t,x)}\mathrm{e}^{\mathrm{i}\theta(0,0)}\big\rangle&=\left(\frac{l^2}{x^2-v^2t^2+\mathrm{i}0^+}\right)^{1/4\pi v\chi}\\
&=\mathrm{e}^{-\mathrm{i}\frac{\pi}{4\pi v\chi}\Theta(v^2t^2-x^2)}\left(\frac{l^2}{|x^2-v^2t^2|}\right)^{1/4\pi v\chi} \tag{3.3.25}
\end{align*}

上述结果告诉我们，在 $1+1$ 维中量子涨落破坏长程序，不存在对称性破缺。但这种破坏并不彻底。$\varphi$ 场仍然具有“准长程序”（quasi-long-range order），其关联按代数方式衰减。这导致如下相图：无论 $\mu<0$ 还是 $\mu>0$，玻色子基态都不破缺 $U(1)$ 对称性。然而，当 $\mu<0$ 时玻色子场只有短程关联；当 $\mu>0$ 时，短程关联被提升为按代数衰减的“准长程”关联。

另一个重要之点是，关联 $\langle\mathrm{e}^{-\mathrm{i}\theta(t,x)}\mathrm{e}^{\mathrm{i}\theta(0,0)}\rangle$ 无法只用有效理论中的参数来表达，它还依赖一个短距离截断标度 $l$。这并不奇怪。要得到低能 XY 模型，我们必须扔掉原理论在短距离和高能量处的大量信息与结构。许多算符的定义本身就依赖这些短距离结构。例如，有些算符是同一空间点上的场量的乘积，而在低能有效理论中，“同一空间点”这一概念已经丧失，至少变得模糊。因此，我们不应指望把所有关联函数都用有效理论中的参数表达。在上面的例子中，我们发现关联 $\langle\mathrm{e}^{-\mathrm{i}\theta(t,x)}\mathrm{e}^{\mathrm{i}\theta(0,0)}\rangle$ 确实依赖理论的短距离结构。令人惊讶的是，所有复杂的短距离结构竟被总结进单一参数 $l$ 之中。这展示了低能有效理论的普适性（universality）。

\subsubsection{$2+1$ 维中的超流相}

为计算 $2+1$ 维的 $G_{\theta}$，我们先计算虚时关联
\begin{equation*}
\mathcal{G}_{\theta}(x^i)=\langle\theta(x^i)\theta(0)\rangle
\end{equation*}
这里 $x^{1,2}$ 是空间坐标，$x^3$ 是虚时间。为简单起见，我们取 $v=1$。XY 模型的配分函数为 $Z=\int\mathcal{D}\theta\,\exp(-\int\mathrm{d}^3x^i\,\frac{\chi}{2}(\partial_i\theta)^2)$。我们看到\footnote{我们用到了 $-\partial_i^2\,r^{-1}=4\pi\delta(x^i)$。}
\begin{equation*}
\mathcal{G}_{\theta}(x^i)=\frac{\chi^{-1}}{-\partial_i^2}=\frac{1}{4\pi\chi r}=\frac{1}{4\pi\chi\sqrt{\pmb{x}^2+\tau^2}}
\end{equation*}
其中 $r^2=x^ix^i$。

由 $\mathrm{i}G_{\theta}$ 与 $\mathcal{G}_{\theta}$ 之间的关系（见式 (2.2.34)），我们得到
\begin{equation*}
\mathrm{i}G_{\theta}(t,x^1,x^2)=\mathcal{G}_{\theta}(x^1,x^2,\mathrm{e}^{\mathrm{i}\pi/2}t)=\frac{1}{4\pi\chi}\frac{1}{\sqrt{\pmb{x}^2-t^2}}
\end{equation*}
然而，当 $|t|>|\pmb{x}|$ 时传播子是含糊的：$\mathrm{i}G_{\theta}(t,x^1,x^2)=\pm\mathrm{i}\frac{1}{4\pi\chi}\frac{1}{\sqrt{|\pmb{x}^2-t^2|}}$。要固定 $\pm$ 号，我们需要更仔细地进行解析延拓。假设 $|t|\gg|\pmb{x}|$，并让 $\mathcal{G}_{\theta}(x^1,x^2,\mathrm{e}^{\mathrm{i}\eta}t)$ 中的 $\eta$ 从 0 连续地变到 $\pi/2$，我们发现
\begin{equation*}
\mathrm{i}G_{\theta}(t,x^1,x^2)=\frac{1}{4\pi v\chi}\frac{1}{\sqrt{|\pmb{x}^2-vt^2|}}\,\mathrm{e}^{-\mathrm{i}\frac{\pi}{2}\Theta(v|t|-|\pmb{x}|)}
\end{equation*}
其中我们已恢复速度 $v$。

无论如何，当 $|t|$ 与 $|\pmb{x}|$ 趋于 $\infty$ 时，$\mathrm{i}G(\pmb{x},t)\to 0$。因此，$\langle\mathrm{e}^{-\mathrm{i}\theta(t,x)}\mathrm{e}^{\mathrm{i}\theta(0,0)}\rangle$ 在长距离处趋于常数 $\mathrm{e}^{-\mathrm{i}G_{\theta}(0,0)}$。然而，该常数的数值不同于经典值（经典值为 1）。引入短距离截断标度 $l$，我们看到
\begin{equation*}
\big\langle\mathrm{e}^{-\mathrm{i}\theta(t,x)}\mathrm{e}^{\mathrm{i}\theta(0,0)}\big\rangle\to\mathrm{e}^{-\frac{1}{4\pi v\chi l}}=\big\langle\mathrm{e}^{-\mathrm{i}\theta(t,x)}\big\rangle\big\langle\mathrm{e}^{\mathrm{i}\theta(0,0)}\big\rangle \tag{3.3.26}
\end{equation*}
我们发现，相涨落在 $2+1$ 维及更高维中不破坏长程序，只是减小序参量的数值。

\subsubsection{短距离截断与非线性效应}

在上面，我们讨论了量子涨落对经典基态的影响。我们发现，在 $1+1$ 维中，量子涨落总是有很大的影响——它们破坏长程序。破坏长程序的量子涨落来自长波长和低频率。无论我们如何调整理论中的参数（例如截断标度），它们都在那里。在 $2+1$ 维及更高维中，长波长、低频涨落不产生发散的效应，一般说来长程序能在量子涨落下存活。序参量敏感地依赖截断标度这一事实表明，序参量的减小量由短距离涨落决定。短距离涨落只是在定量上修正经典结果。

现在产生了一个问题：这些定量修正何时很小？也就是说，经典近似何时能定量地描写系统的物理性质？从关于序参量的结果（见式 (3.3.26)）我们看到，如果截断标度 $l\gg(v\chi)^{-1}$，涨落修正就很小。为得到截断标度，注意我们在式 (3.3.9) 中略去了 $\frac{\varphi_0^2}{2m}\partial_{\pmb{x}}^2$ 项，因为我们假定 $\theta$ 的涨落是平滑的。然而，在短距离标度处，当 $\frac{\varphi_0^2}{2m}\partial_{\pmb{x}}^2$ 与另一项 $2V_0\varphi_0^4$ 可比时，我们不能再忽略梯度项，XY 模型失效。这一跨越长度标度（crossover length scale）由
\begin{equation*}
\xi=(4mV_0\varphi_0^2)^{-1/2}
\end{equation*}
给出，称为相干长度（coherence length）。相干长度有如下含义：如果我们改变某一点的 $h$ 场或玻色子密度，那么这种改变将在 $\xi$ 的距离上传播。取 $l=\xi$，约化因子变为
\begin{equation*}
\mathrm{e}^{-\frac{1}{4\pi v\chi l}}=\mathrm{e}^{-\frac{mV_0}{\sqrt{2}\pi}} \tag{3.3.27}
\end{equation*}
我们可以把上述结果改写成一个更有意义的形式。注意 $E_{int}\equiv\frac{1}{2}V_0\varphi_0^2$ 具有能量的量纲，它表示每个粒子的相互作用能。在 $2+1$ 维中，$E_{qua}\equiv\rho_0/m$ 也具有能量的量纲。$E_{qua}$ 是这样一个能标（或温度标度），低于它时玻色子波函数开始交叠，我们需要把玻色子当作量子系统来处理。现在我们可以把约化因子改写为
\begin{equation*}
\mathrm{e}^{-\frac{1}{4\pi v\chi l}}=\mathrm{e}^{-\frac{\sqrt{2}E_{int}}{E_{qua}}}
\end{equation*}
我们看到，在弱相互作用极限下涨落修正很小，经典结果良好。

然而，上述结果并不完整，因为我们只在二次近似之内考虑了短距离涨落的效应，还没有考虑高阶项带来的非线性效应。为了更系统地研究涨落的效应，我们想对我们的玻色子作用量作量纲分析（dimensional analysis）：
\begin{equation*}
S=\int\mathrm{d}^d\pmb{x}\,\mathrm{d}t\,\Big[\mathrm{i}\frac{1}{2}\big(\varphi^*\partial_t\varphi-\varphi\partial_t\varphi^*\big)-\frac{1}{2m}\partial_{\pmb{x}}\varphi^*\partial_{\pmb{x}}\varphi-\frac{V_0}{2}|\varphi|^2\big(|\varphi|^2-2\rho_0\big)\Big] \tag{3.3.28}
\end{equation*}
这里我们取了 $\mu=V_0\rho_0$。我们想重新标度 $t$、$\pmb{x}$ 和 $\varphi$，把作用量改写成 $S=\tilde{S}/g$ 的形式，使得 $\tilde{S}$ 中所有的系数都是 1 的量级。我们发现下述重新标度即可奏效：
\begin{equation*}
x=\xi\tilde{x},\qquad t=(V_0\rho_0)^{-1}\tilde{t},\qquad\varphi=\sqrt{\rho_0}\,\tilde{\varphi}
\end{equation*}
它给出
\begin{equation*}
S=g^{-1}\int\mathrm{d}^d\tilde{\pmb{x}}\,\mathrm{d}\tilde{t}\,\Big[\mathrm{i}\frac{1}{2}\big(\tilde{\varphi}^*\partial_t\tilde{\varphi}-\tilde{\varphi}\partial_t\tilde{\varphi}^*\big)-2\partial_{\pmb{x}}\tilde{\varphi}^*\partial_{\pmb{x}}\tilde{\varphi}-\frac{1}{2}|\tilde{\varphi}|^2\big(|\tilde{\varphi}|^2-2\big)\Big]
\end{equation*}
其中
\begin{equation*}
g=N_{\xi}^{-1},\qquad N_{\xi}=\rho_0\xi^d
\end{equation*}
这里 $N_{\xi}$ 是体积 $\xi^d$ 中的粒子数。我们也可以把 $g$ 写成
\begin{equation*}
g=\sqrt{\rho_0^{d-2}(4mV_0)^d} \tag{3.3.29}
\end{equation*}

现在很清楚，如果路径积分 $\int\mathcal{D}^2\tilde{\varphi}\,\mathrm{e}^{\mathrm{i}\frac{\tilde{S}}{g}}$ 中的 $g$ 很小，那么“势”很陡，势能极小值附近的涨落很小。这时，半经典近似（semiclassical approximation）是好的。实际上，半经典近似正对应于按 $g$ 的展开。在 $2+1$ 维中 $g=4mV_0$，它恰好与从短距离涨落得到的条件 (3.3.27) 相符。

当 $g$ 大时，涨落可以大到这样的程度，以至于经典图像在定性上是否正确都不清楚了。事实上，人们相信，当 $g\gg1$ 时，短距离涨落可以破坏长程序并在基态中恢复 $U(1)$ 对称性，例如通过形成晶体（晶体具有不同的对称性破缺和不同的长程序）。这与 $1+1$ 维的情形形成对照，在那里无论 $g$ 取何值，长距离涨落都破坏长程序。

按照经典理论，我们的玻色子系统在 $\mu=0$ 处经历一个量子相变（quantum phase transition）。该相变是连续的，由一个量子临界点描写。在经典理论之内，我们可以计算出所有的临界指数。例如，$\omega\propto k^z$ 中的动力学指数（dynamical exponent）$z$ 为 $z=2$；$\langle\varphi\rangle=\varphi_0\propto\mu^{\nu}$ 中的指数 $\nu$ 为 $\nu=1/2$。问题是，经典理论能正确地描写量子临界点吗？由式 (3.3.29) 可见，如果 $d>2$，那么我们离临界点（$\rho_0\to0$）越近，$g$ 越小，半经典近似就越好。因此，对于 $d>2$，经典理论能正确地描写量子临界点。然而，对于 $d<2$，当我们逼近临界点时 $g$ 发散。在这种情况下，我们必须把量子涨落包括进来才能得到正确的临界指数。这一跨越空间维数 $d_c=2$ 称为该临界点的上临界维数。

\subsection*{问题 3.3.5}

考虑一个具有长程相互作用的玻色子系统
\begin{equation*}
\frac{1}{2}\int\mathrm{d}^d\pmb{x}\,\mathrm{d}^d\pmb{x}'\,|\varphi(\pmb{x})|^2V(\pmb{x}-\pmb{x}')|\varphi(\pmb{x}')|^2
\end{equation*}
其中 $V(r)=V_dr^{\epsilon-d}$。如果 $\epsilon<0$，则相互作用实际上是短程的。这里我们假设 $\epsilon>0$。
\begin{enumerate}
\item 推导修改后的 XY 模型。
\item 找出临界空间维数 $d_l$，低于它时相位涨落 $\theta$ 总是破坏长程序。这里我们把维数当作连续的实数处理。我们知道，对于短程相互作用 $d_l=1$。
\item 重复本节末尾的讨论，求出 $g$（它决定量子涨落何时很小）以及上临界维数 $d_c$（超过它时经典理论能正确描写量子临界点）。
\end{enumerate}

\subsection*{问题 3.3.6}

$1+1$ 维中的有限温关联。
\begin{enumerate}
\item 计算虚时关联
\begin{equation*}
\mathcal{G}(x,\tau)=\big\langle\mathrm{e}^{\mathrm{i}\theta(x,\tau)}\mathrm{e}^{-\mathrm{i}\theta(0,0)}\big\rangle
\end{equation*}
假设一维空间是长度为 $L$ 的有限圆。
\item 计算有限温下的虚时关联 $\mathcal{G}^{\beta}(x,\tau)$，假设一维空间是一条无限长的直线。（提示：你可以交换 $x$ 与 $\tau$，并利用第 1 问的结果。）
\item 计算有限温下的实时（编时）关联 $G^{\beta}(x,t)$，假设一维空间是一条无限长的直线。证明
\begin{equation*}
G^{\beta}(0,t)\propto\mathrm{e}^{-\mathrm{i}\frac{\pi}{4\pi v\chi}}\left(\frac{\pi T}{\sinh(\pi T|t|)}\right)^{1/2\pi v\chi}
\end{equation*}
\end{enumerate}

\section{有限温度下的超流相}

\subsection{有限温度下的路径积分}

\begin{keypoints}
\item $(d+1)$ 维量子系统在有限温度下的长距离性质，可以用一个 $d$ 维系统的路径积分来描写。
\end{keypoints}

本节中，我们将研究一个相互作用玻色子系统及其有限温度下的超流相。我们从虚时路径积分
\begin{equation*}
Z=\int\mathcal{D}^2[\varphi(\pmb{x},\tau)]\,\mathrm{e}^{-\int_0^{\beta}\mathrm{d}^d\pmb{x}\,\mathrm{d}\tau\,\big[\frac{1}{2}(\varphi^*\partial_{\tau}\varphi-\varphi\partial_{\tau}\varphi^*)+\frac{1}{2m}\partial_{\pmb{x}}\varphi^*\partial_{\pmb{x}}\varphi-\mu|\varphi|^2+\frac{V_0}{2}|\varphi|^4\big]}
\end{equation*}
出发，它表示系统的配分函数。注意，贝里相项在虚时路径积分中保持为虚的。

为了把上述配分函数约化为一个熟悉的统计模型，我们先通过引入
\begin{equation*}
\varphi_{\omega_n}(\pmb{x})=\beta^{-1}\int_0^{\beta}\mathrm{d}\tau\,\varphi(\pmb{x},\tau)\mathrm{e}^{\mathrm{i}\tau\omega_n}
\end{equation*}
进入离散频率空间。然后，我们积掉所有有限频率的模式：
\begin{equation*}
Z=\int\prod_n\mathcal{D}^2[\varphi_{\omega_n}]\,\mathrm{e}^{-S}=\int\mathcal{D}^2[\varphi_c(\pmb{x})]\,\mathrm{e}^{-S_{\mathrm{eff}}}
\end{equation*}
其中 $\varphi_c(\pmb{x})=\varphi_{\omega_n=0}(\pmb{x})$ 是零频率模式。这一步很困难，所得的有效作用量 $S_{\mathrm{eff}}$ 可以非常复杂，甚至可以包含非局域项，例如 $\int\mathrm{d}^d\pmb{x}\,\mathrm{d}^d\pmb{x}'\,|\varphi_c(\pmb{x})|^2K(\pmb{x}-\pmb{x}')|\varphi_c(\pmb{x}')|^2$。然而，在有限频率处，原作用量包含 $\mathrm{i}|\varphi_c(\pmb{x})|^2\omega_n$ 项，它使 $\varphi_n$ 的传播子具有 $1/(\mathrm{i}\omega_n+ck^2+\mu)$ 的形式。该传播子是短程的，按指数衰减。因此，非零模式的涨落只能传递短程相互作用，$K(\pmb{x}-\pmb{x}')$ 应具有指数衰减：$K(\pmb{x}-\pmb{x}')\sim\mathrm{e}^{-|\pmb{x}-\pmb{x}'|/l_T}$。这里我们有一个新的长度标度 $l_T$，超出它之后，有效作用量 $S_{\mathrm{eff}}$ 就可以当作局域作用量处理。尽管 $S_{\mathrm{eff}}$ 可以非常复杂，但由于其对称性，局域作用量只能取某种确定的形式。在超出 $l_T$ 的长距离处，我们可以作梯度展开（gradient expansion）
\begin{equation*}
S_{\mathrm{eff}}=\beta\int\mathrm{d}^dx\,\big[\frac{1}{2m^*}|\partial_x\varphi_c|^2+V(|\varphi_c|,T)+\dots\big] \tag{3.4.1}
\end{equation*}
“…”代表更高阶导数项。有效势（effective potential）依赖于温度。当 $\mu<0$ 时，我们预期 $V$ 总有一个位于 $\varphi_c=0$ 的单一极小，代表对称态。当 $\mu>0$ 且温度较低时，$V$ 应有一个位于 $\varphi_c=\varphi_0\mathrm{e}^{\mathrm{i}\theta}$ 的极小圆圈，代表对称性破缺态。然而，超出临界温度 $T_c$ 之后，预期 $V$ 会变成在 $\varphi_c=0$ 处只有一个单一极小的形式。因此，$V$ 对温度的依赖描写了超流转变。

% ==================================================================
% 新增术语（ch3_c 新增，供术语表追加）：
%   canonical momentum 正则动量；uncertainty relation 不确定关系；
%   on-site repulsion 在位排斥；lattice boson system 格点玻色子系统；
%   compressibility 压缩率；Nambu–Goldstone modes Nambu–Goldstone 模；
%   quasi-long-range order 准长程序；universality 普适性；
%   short-distance cut-off (scale) 短距离截断（标度）；crossover length scale 跨越长度标度；
%   coherence length 相干长度；interaction energy per particle 每粒子相互作用能；
%   dimensional analysis 量纲分析；semiclassical approximation 半经典近似；
%   quantum phase transition 量子相变；quantum critical point 量子临界点；
%   dynamical exponent 动力学指数；upper critical dimension 上临界维数；
%   quantum critical exponents 量子临界指数；long-range interaction 长程相互作用；
%   zero-frequency mode 零频率模式；discrete frequency space 离散频率空间；
%   non-local term 非局域项；gradient expansion 梯度展开；
%   effective potential 有效势；phase fluctuation 相位涨落
% ==================================================================
